% id: 143-20
% book: 베이직쎈 대수 · 08 삼각함수의 활용
% section: 유형 084 사인법칙과 코사인법칙의 실생활에의 활용
% figure: tikz:fig-143-20
% answer: 
% answer_source: 
% variant_level: 0 (원본 전사)
% 이 파일은 build.mjs 가 items.json 에서 만든다. 고칠 때는 items.json 을 고친다.
\smash{\raisebox{4.5mm}{\dmkichul{베이직쎈 대수 143쪽 20번}}}
\noindent\begin{minipage}[t]{0.58\linewidth}\vspace{0pt}\raggedright 오른쪽 그림과 같이 $50\sqrt{3}\mkern3mu \mathrm{m}$ 떨어진 두 지점~$\pt{A}$, $\pt{B}$에서 분수대가 있는 $\pt{C}$ 지점을 바라본 각의 크기를 재었더니 $\angle\pt{A}=75^\circ$, $\angle\pt{B}=45^\circ$이었다. 세 지점~$\pt{A}$, $\pt{B}$, $\pt{C}$를 연결하는 원 모양의 잔디밭을 만들려고 할 때, 잔디밭의 넓이를 구하시오.\end{minipage}\hfill\begin{minipage}[t]{0.4\linewidth}\vspace{0pt}\centering \resizebox{\ifdim\width>\linewidth\linewidth\else\width\fi}{!}{% 143-20(143쪽) — 세 지점 A, B, C 를 지나는 원 모양의 잔디밭과 삼각형 ABC · ∠A=75°, ∠B=45°, AB=50√3 m.
% 스캔 삽화(연두색 잔디밭 바탕·분수대 그림)는 화질이 낮아 TikZ 로 다시 그렸다 — 원(잔디밭)은 문제의 조건이라 남기고 분수대 그림은 뺐다.
% 라벨은 원문에 있는 A·B·C·75°·45°·50√3 m 만 두고 잔디밭의 넓이(답)는 적지 않는다.
\begin{tikzpicture}[scale=1.15, line width=0.5pt, >=stealth]
  \coordinate (A) at (100:1);
  \coordinate (B) at (220:1);
  \coordinate (C) at (10:1);
  \fill[yellow!35] (0,0) circle (1);
  \draw[line width=0.4pt, yellow!60!black] (0,0) circle (1);
  \draw (A) -- (B) -- (C) -- cycle;
  \draw[line width=0.4pt] ([shift=(250:0.24)]A) arc (250:325:0.24);
  \draw[line width=0.4pt] ([shift=(25:0.26)]B) arc (25:70:0.26);
  \node[font=\footnotesize] at ([shift=(288:0.4)]A) {$75^\circ$};
  \node[font=\footnotesize] at ([shift=(48:0.44)]B) {$45^\circ$};
  \node[above=0pt, font=\small] at (A) {$\pt{A}$};
  \node[below=0pt, font=\small] at (B) {$\pt{B}$};
  \node[below right=-3pt and -2pt, font=\small] at (C) {$\pt{C}$};
  \draw[densely dashed, line width=0.3pt, <->] (-0.400,1.077) -- (-0.992,-0.550);
  \node[left=2pt, font=\footnotesize] at (-0.700,0.225) {$50\sqrt{3}\,\mathrm{m}$};
\end{tikzpicture}
}\end{minipage}\par
